\documentclass[9pt, french]{article}
\usepackage{graphicx}
\usepackage[T1]{fontenc}
\usepackage[utf8]{inputenc}
\usepackage{lmodern}
\usepackage[a4paper]{geometry}
\usepackage{babel}
\usepackage{wrapfig}
\usepackage[figurename=Fig.]{caption}
\usepackage{subfig}
\begin{document}
	\begin{titlepage}
		
		\newcommand{\HRule}{\rule{\linewidth}{0.5mm}} % Defines a new command for the horizontal lines, change thickness here
		
		\center % Center everything on the page
		
		%----------------------------------------------------------------------------------------
		%	HEADING SECTIONS
		%----------------------------------------------------------------------------------------
		
		\textsc{\LARGE Université de Montréal}\\[1cm] % Name of your university/college
		\textsc{\Large Faculté des arts et des sciences}\\[0.5cm] % Major heading such as course name
		\textsc{\large Département d’informatique et de recherche opérationnelle (DIRO) }\\[0.5cm] % Minor heading such as course title
		
		%----------------------------------------------------------------------------------------
		%	TITLE SECTION
		%----------------------------------------------------------------------------------------
		
		\HRule \\[0.4cm]
		{ \huge Rapport de stage \\ \bfseries Évaluation du risque de retour à la maison } \\[0.4cm] % Title of your document
		\HRule \\[1.5cm]
		
		%----------------------------------------------------------------------------------------
		%	AUTHOR SECTION
		%----------------------------------------------------------------------------------------
		
		\begin{minipage}{0.4\textwidth}
			\begin{flushleft} \large
				\emph{Auteur:}\\
				Vilon \textsc{Saint-Fleurose} % Your name
				\\MSc, informatique \\Université de Montréal
			\end{flushleft}
		\end{minipage}
		~
		\begin{minipage}{0.4\textwidth}
			\begin{flushright} \large
				\emph{Directeur de recherche:} \\
				Dr. Michalis \textsc{Famelis} % Supervisor's Name
				\\Professeur adjoint \\Université de Montréal
			\end{flushright}
		\end{minipage}\\[2cm]
			~
			\begin{minipage}{0.4\textwidth}
				\begin{center} \large
					\emph{Superviseur:} \\
					 Nicolas \textsc{Coallier} % Supervisor's Name
					\\Vice-Président Exécutif, TIC \\ML+
				\end{center}
			\end{minipage}\\[2cm]
		
		% If you don't want a supervisor, uncomment the two lines below and remove the section above
		%\Large \emph{Author:}\\
		%John \textsc{Smith}\\[3cm] % Your name
		
		%----------------------------------------------------------------------------------------
		%	DATE SECTION
		%----------------------------------------------------------------------------------------
		
		{\large \today}\\[2cm] % Date, change the \today to a set date if you want to be precise
		
		%----------------------------------------------------------------------------------------
		%	LOGO SECTION
		%----------------------------------------------------------------------------------------
		
		\includegraphics[width=4cm, height=1.2cm]{logo.png} % Include a department/university logo - this will require the graphicx package
		
		%----------------------------------------------------------------------------------------
		
		\vfill % Fill the rest of the page with whitespace
		
	\end{titlepage}
	
	\section*{Remerciements}
	Je veux commencer d'abord par remercier le Grand Dieu Tout-Puissant, le Créateur de l'univers, des cieux et de la terre qui m'a donné la vie, la santé, les opportunités et tout ce dont j'avais besoin pour faire cette grande et belle étude à l'université de Montréal. Il a rendu toutes choses possibles en ma faveur, moi qui suis pécheur et désobéissant; immérité de toutes ces grâces. Il m'accompagnait toujours dans les moments les plus difficiles de ma vie, Il ne m'a jamais laissé seul; surtout dans les moments où je devais payer les frais de scolarité qui étaient si énormes et impossibles à payer de mon propre compte. A un Dieu si merveilleux et si bon, je Lui dois beaucoup de reconnaissance.  
	
	Je remercie aussi ma femme qui m'a beaucoup supporté pendant plus de deux années d'études. Elle n'a jamais murmuré, ni découragé quand nous devions passer par des moments difficiles de notre vie conjugale à cause de ces études. Elle a mis toutes ses ressources disponibles pour entretenir la famille et payer mes études quand j’étais moi-même dans l’impossibilité de travailler. Vraiment, ma femme est une bénédiction dans ma vie, un cadeau venant de Dieu. Je t’aime ma chérie.
	
	Je tiens à remercier le professeur Michalis Famelis d’avoir accepté être mon directeur de recherche et supervisé ce stage, il est toujours là pour m’encourager et me pousser vers l’avant. Il répond toujours présent à tous mes appels, il est toujours disponible pour me rencontrer, me parler et me conseiller ; même en dehors du cadre universitaire. 
	
	Je remercie Nicolas Coallier et toute l’équipe ML+ qui ont accepté que je sois leur stagiaire, ils ont placé  leur confiance en moi quoiqu’ils ne me connaissaient pas encore. Cette équipe, quoique jeune, est très dynamique et chaleureuse, c’est une équipe motivante qui stimule la connaissance. J’ai dû apprendre beaucoup de choses par rapport à eux.
	Finalement, je présente mes sincères remerciements à toute la communauté universitaire, à DIRO en particulier. Merci pour la formation prestigieuse que vous m’avez fournie. Cette formation est si solide qu’elle m’aidera rapidement à intégrer le marché du travail sans perdre de temps. 
	
	\newpage
	

	\section{Introduction}
	
	Le système de santé aux USA repose sur deux type de financement. Premièrement, sur le financement public qui englobe certains groupes de personnes de la population américiane uniquement. Il se concentre essentiellement sur deux programmes : le programme 		fédéral Medicare pour les plus de 65 ans et les personnes gravement handicapées (soit 15 \% de la population) et le programme Medicaid qui s’adresse aux familles pauvres avec enfants et touche 11 \% de la population. Deuxièmement, sur le financement privé qui touche tout le reste de la population, l’assurance est donc majoritairement privée aux Etats-Unis. Les Américains sont assurés en général via leurs employeurs ou sinon de manière individuelle lorsque leur employeur ne propose pas d’assurance ou qu’ils travaillent en indépendant. La composante « assurance médicale » dans le choix d’un emploi est donc un critère important.
	
	La science des données est un domaine interdisciplinaire de méthodes, processus, algorithmes et systèmes scientifiques permettant d'extraire des connaissances ou des informations à partir de données sous diverses formes, structurées ou non. Elle emploie des techniques et des théories tirées de plusieurs autres domaines plus larges des mathématiques, la statistique principalement, la théorie de l'information et la technologie de l'information, notamment le traitement de signal, des modèles probabilistes, l'apprentissage automatique, l'apprentissage statistique, la programmation informatique, l'ingénierie de données, la reconnaissance de formes et l'apprentissage, la visualisation, l'analytique prophétique, la modélisation d'incertitude, le stockage de données, la compression de données et le calcul à haute performance. Les méthodes qui s'adaptent aux données de masse sont particulièrement intéressantes dans la science des données, bien que la discipline ne soit généralement pas considérée comme limitée à ces données.\\
	
		

	\section{Solutions techniques proposées}
	
	

	
	
\end{document}
